\newif\ifglobal

%%%%%%%%%%%%%%%%%%%%%%%
%%% Compile to view %%%
%%%%%%%%%%%%%%%%%%%%%%%

%%%%%%%%%%%%%%%%%%%%%%%%%%%%%%%%%%%%%%%%%%%%%%%%%%%
%%% Readme:                                     %%%
%%% https://www.overleaf.com/read/bwdjvfjksvjy/ %%%
%%%%%%%%%%%%%%%%%%%%%%%%%%%%%%%%%%%%%%%%%%%%%%%%%%%

% >>> M?: marks a module setting number or important notice

% >>> M4: Set {./relative-path} to external files for this module <<<
\def \modulefiles {../ESP32-C61/23-ECDSA}
% To add an external file, use:
% (Replace '---' with actual filename)
% '\input{\modulefiles/tables/---.tex}' for tables
% '\includegraphics{\modulefiles/figures/---}' for figures
% '\subfile{\modulefiles/---}' for register files

% >>> M1: This module is in global project? <<<
% 'YES' - uncomment; 'NO' - comment out
\globaltrue

\ifglobal
    \documentclass[main\_\_EN.tex]{subfiles}
   
\else
    \documentclass[openany, 10pt]{book}

    % >>> M2: In ./00-shared/preamble-trm-module, also find
    %         settings for visibility of todo notes, watermark <<<
    \usepackage{preamble-trm-module}

    % >>> M3: Temporary labels in English,
    %         you can add your own labels there <<<
    \myexternaldocument{./00-shared/temp-labels-en}

\fi %\ifglobal


\setcounter{secnumdepth}{4}


% Define formatting of parts and chapters in text
\titleformat{\chapter}
  [display]
  {\LARGE\bfseries}
  {\color{AllportsColor}Chapter \thechapter}
  {1em}
  {\color{AllportsColor}}
\titlespacing*{\chapter}{0pt}{0pt}{20pt}

% Define headers
\usepackage{fancyhdr}
\usepackage{lipsum}
\renewcommand{\chaptermark}[1]{\markboth{Chapter\, \thechapter\, \ #1}{}}
\fancypagestyle{plain}{%
  \fancyhead{}%
  %\fancyhead[R]{\Navigation}
  \fancyhead[L]{%
    \nouppercase{%
      \ %
      \normalsize\leftmark
    }%
  }%
}
\pagestyle{plain}

% >>> M5: If this module needs any updates in
% './00-shared/preamble-trm-module.sty'
% (include additional package, etc.), contact your module PM
% or create the file './00-shared/changelog.tex' make the
% necessary updates yourself and document them in this file <<<

% Make text left-aligned and allow latex to break pages naturally, comment out for full justification
\raggedright
\raggedbottom


\begin{document}


% Sets language into which pre-defined bits of text are translated
\selectlanguage{english}

% Do not delete this block
\ifglobal
% Add the following front matter if
% this module is outside of global project
\else

    \InputIfFileExists{readme}{}{}
    \listoftodos
    {\let\clearpage\relax \listofdonetodos}
    \tableofcontents
    %\listoftables
    %\listoffigures
\fi


%%%%%%%%%%%%%%%%%%%%%%%%%
%%% TRM Module Begins %%%
%%%%%%%%%%%%%%%%%%%%%%%%%

\hypertarget{ecdsa}{}
\chapter{ECDSA Digital Signature Peripheral (ECDSA\_DS)}
\label{mod:ecdsa}

%%%%%%%%%%%%%%%%%%%%
%%% Introduction %%%
%%%%%%%%%%%%%%%%%%%%

\section{Introduction}

In cryptography, the Elliptic Curve Digital Signature Algorithm (ECDSA) offers a variant of the Digital Signature Algorithm (DSA) which uses elliptic-curve cryptography.

\chipname{}'s ECDSA Digital Signature Peripheral (ECDSA\_DS) provides a secure and efficient environment for computing ECDSA signatures. It offers fast computations while ensuring the confidentiality of the signing process to prevent information leakage. This makes it a valuable tool for applications that require high-speed cryptographic operations with strong security guarantees. By using ECDSA\_DS, users can be confident that their data is protected without sacrificing performance.


%%%%%%%%%%%%%%%%%%%%
%%% Feature List %%%
%%%%%%%%%%%%%%%%%%%%

\section{Features}
\label{sec:ecdsa-features}

\chipname{}'s ECDSA\_DS module supports:

\begin{itemize}
    \item Digital signature generation and verification
    \item Two different elliptic curves, namely P-192 and P-256, defined in \href{https://nvlpubs.nist.gov/nistpubs/FIPS/NIST.FIPS.186-5.pdf}{FIPS 186-5}
    \item Two hash algorithms for message hash in the ECDSA\_DS operation, namely SHA-224 and SHA-256, defined in \href{https://doi.org/10.6028/NIST.FIPS.180-4}{FIPS PUB 180-4 Spec}
    \item State-dependent register access control in different operation states to ensure information security
\end{itemize}


%%%%%%%%%%%%%%%%%%%
%%% Terminology %%%
%%%%%%%%%%%%%%%%%%%

\label{sec:ecdsa-terminology}

\section{ECDSA Basics}
\subsection{Domain Parameters}

ECDSA uses parameters that define the elliptic curve over a finite field, as well as the generator point and the order of the base point. These parameters are usually referred to as domain parameters, and they are required for key generation, signature generation, and signature verification.

The domain parameters used in ECDSA consist of the following:

\begin{itemize}
    \item The elliptic curve domain parameters, which include:
            \begin{itemize}
                \item The prime modulus $p$, which specifies the size of the finite field over which the elliptic curve is defined.
                \item The coefficients $a$ and $b$, which specify the elliptic curve equation.
                \item The base point $G$ on the curve, which can be used to generate public keys.
                \item The order $n$ of the base point, which is the number of points on the curve that can be generated by repeatedly adding $G$ to itself.
            \end{itemize}
    \item The parameters for the hash function, which include:
          \begin{itemize}
              \item The hash algorithm used to generate a fixed-length hash value from the message being signed.
              \item The length of the hash value, which determines the size of the signature.
          \end{itemize}
\end{itemize}

Together, these parameters define the ECDSA domain.

\subsection{Key Generation}

The process of key generation in ECDSA is described below:
\begin{enumerate}
    \item Select an elliptic curve domain by defining:
    \begin{itemize}
        \item the prime modulus $p$
        \item the coefficients $a$ and $b$
        \item the base point $G$ and its order $n$
    \end{itemize}
    \item Generate a private key:
    \begin{itemize}
        \item A private key $d$ is a randomly chosen integer between 1 and $n$-1. It is essential to use a secure \textit{\nameref{mod:rng}} to ensure that the private key is truly random and cannot be predicted or reproduced.
        \item The private key is used for signature generation. This key must be kept secret and secure.
        \item Once a private key is chosen, it should be burned into the eFuse in \chipname{}. (Please refer to Chapter \ref{mod:efuse} \textit{\nameref{mod:efuse}}).
    \end{itemize}
    \item Compute the public key:
    \begin{itemize}
        \item The public key $Q$ is calculated as $Q = dG$.
    \end{itemize}
    \item Export the public key:
    \begin{itemize}
        \item The public key $Q$ is typically represented as a pair of coordinates ($Qx$, $Qy$) that can be shared with others.
    \end{itemize}
\end{enumerate}

\subsection{Signature Generation}

The ECDSA signature generation process is described below:

\begin{enumerate}
    \item Select a message $m$ to be signed.
    \item Calculate the hash of the message $e$: $e$ is equal to HASH(m), where HASH is a cryptographic hash function, such as SHA-256.
    \item Compute the digest of the message hash $z$: Let $z$ be the $L_n$ leftmost bits of $e$, where $L_n$ is the bit length of the base point order $n$.
    \item Select a random number $k$, which is chosen between 1 and $n$-1, where $n$ is the order of the base point on the elliptic curve.
    \item Compute the signature: The signature is calculated as follows:
          \begin{enumerate}
              \item Compute the point ($x$, $y$) = $kG$
              \item Calculate $r$ = $x$ mod $n$. If $r$ is equal to zero, return to the previous step and select a new random value of $k$.
              \item Calculate $s$ = $k$$^{-1}$ * $(z + d*r)$  \textrm{mod} $n$. If $s$ is equal to zero, return to step 4 and select a new random value of $k$.
              \item The signature is the pair ($r$, $s$).
          \end{enumerate}
    \item Send the message $m$ and signature ($r$, $s$) to the recipient.
\end{enumerate}

\subsection{Signature Verification}

The recipient can then use the public key associated with the private key used for signature generation to verify the signature. The verification process involves checking that the signature was generated using the correct private key and that the signature is valid for the given message.

The ECDSA signature verification process is described below:

\begin{enumerate}
    \item Obtain public key $Q$ and receive the message $m$ and signature ($r$, $s$).
    \item Calculate the hash of the message $e$: $e$ is equal to HASH(m), where HASH is a cryptographic hash function, such as SHA-256.
    \item Compute the digest of the message $z$: Let $z$ be the $L_n$ leftmost bits of $e$, where $L_n$ is the bit length of the base point order $n$.
    \item Verify the signature: The signature is verified as follows:
        \begin{enumerate}
            \item Verify that $r$ and $s$ are integers between 1 and $n$-1, where $n$ is the order of the base point on the elliptic curve. If either $r$ or $s$ is outside of this range, the signature is invalid.
            \item Calculate $u_1 = z * s^{-1}$ mod $n$ and $u_2 = r * s^{-1}$ mod $n$.
            \item Calculate the point $(x_1, y_1) = u_1 * G + u_2 * Q$, where $G$ is the base point on the elliptic curve, and $Q$ is the public key associated with the private key used for signature generation.
            \item Verify that $r = x_1$ \textrm{mod} $n$. If $r$ is not equal to $x_1$ mod $n$, the signature is invalid.
        \end{enumerate}
    \item Accept or reject the signature: If the signature is valid, the recipient can be confident that the message was not tampered with and that it came from the expected sender, thus can accept the message as authentic. Otherwise, the recipient rejects the message as invalid.
\end{enumerate}

\section{Functional Description}

This section describes the details of \chipname{}'s ECDSA\_DS module.

\subsection{ECDSA\_DS Working Modes}

The ECDSA\_DS module integrated in the \chipname{} has three working modes, which are Signature Generation, Signature Verification and Public Key Export modes.

Users can select the working mode for ECDSA\_DS by configuring \hyperref[fielddesc:ECDSAWORKMODE]{ECDSA\_WORK\_MODE} according to Table \ref{tab:ecdsa-working-mode} below.

\begin{table}[H]
    \centering
    \caption{ECDSA\_DS Working Mode}
    \label{tab:ecdsa-working-mode}
    \begin{tabular}{ | c | l | }
    \hline
        \rowcolor{lightgray}
        \textbf{\hyperref[fielddesc:ECDSAWORKMODE]{ECDSA\_WORK\_MODE}}      & \textbf{Working Mode}\\\hline
        0                       & Signature Verification    \\\hline
        1                       & Signature Generation      \\\hline
        2                       & Public Key Export         \\\hline
    \end{tabular}
\end{table}

Users can select the elliptic curves used by the ECDSA\_DS module by configuring \hyperref[fielddesc:ECDSAECCCURVE]{ECDSA\_ECC\_CURVE} according to Table \ref{tab:ecdsa-curve-selection}.

\begin{table}[H]
    \centering
    \caption{ECDSA\_DS Elliptic Curve Selection}
    \label{tab:ecdsa-curve-selection}
    \begin{tabular}{ | c | l | }
    \hline
        \rowcolor{lightgray}
        \textbf{\hyperref[fielddesc:ECDSAECCCURVE]{ECDSA\_ECC\_CURVE}}      & \textbf{Elliptic Curve}\\
        \hline
        0                       & P-192    \\\hline
        1                       & P-256    \\\hline
    \end{tabular}
\end{table}

Users can select the SHA algorithms for message hash by configuring \hyperref[fielddesc:ECDSASHAMODE]{ECDSA\_SHA\_MODE} according to Table \ref{tab:ecdsa-sha-algorithms}.

\begin{table}[H]
    \centering
    \caption{ECDSA\_DS SHA Algorithm}
    \label{tab:ecdsa-sha-algorithms}
    \begin{tabular}{ | c | l | }
    \hline
        \rowcolor{lightgray}
        \textbf{\hyperref[fielddesc:ECDSASHAMODE]{ECDSA\_SHA\_MODE}}      & \textbf{SHA Algorithm}\\
        \hline
        1                       & SHA-224    \\\hline
        2                       & SHA-256    \\\hline
        Others                  & Invalid    \\\hline
    \end{tabular}
\end{table}

Additionally, users can check the working states of ECDSA\_DS by inquiring the \hyperref[regdesc:ECDSASTATEREG]{ECDSA\_STATE\_REG} register and comparing the return value against the Table \ref{tab:ecdsa-working-state} below.

\begin{table}[H]
    \centering
    \caption{ECDSA\_DS Working State}
    \label{tab:ecdsa-working-state}
    \begin{tabular}{ | c | l | L{8cm} | }
    \hline
        \rowcolor{lightgray}
        \textbf{\hyperref[regdesc:ECDSASTATEREG]{ECDSA\_STATE\_REG}} & \textbf{State} &  \textbf{Description} \\ \hline
        0    & IDLE      &   Idle or completed operation. Corresponding to IDLE stage.\\ \hline
        1    & LOAD      &   Waiting for users to load information into ECDSA\_DS. Corresponding to LOAD stage.    \\ \hline
        2    & GAIN      &   Waiting for users to gain information from ECDSA\_DS. Corresponding to GAIN stage.    \\ \hline
        3    & BUSY      &   In the middle of a hardware operation. Corresponding to PREP, PROC, and POST stage.    \\ \hline
    \end{tabular}
\end{table}

\subsection{Data and Data Block}

\chipname{}'s ECDSA\_DS module operates on data of 256 or 512 bits. This data ($D[255:0]$) can be divided into 32-bit data blocks.\newline
Take 256-bit long data as an example, $D[n][31:0] (n=0, 1, \cdots, 7)$. Data blocks with the smaller serial number correspond to the lower binary bits. To be specific:
\[
D[255:0] = {D[7][31:0],D[6][31:0],D[5][31:0],
D[4][31:0],D[3][31:0],D[2][31:0],D[1][31:0],D[0][31:0]}
\]

\subsubsection{Writing Data}

Writing data means writing data to an ECDSA\_DS memory block and using this data as the input to ECDSA\_DS. To be specific, writing data to an ECDSA\_DS memory block means writing $D[n][31:0]$ to the ``starting address of this ECDSA\_DS memory block + $4 \times n$".
For a 256-bit long data example:
\begin{itemize}
    \item write $D[0]$ to ``starting address"
    \item write $D[1]$ to ``starting address + 4"
    \item $\cdots$
    \item write $D[7]$ to ``starting address + 28"
\end{itemize}

\begin{tiplisting}
When the data size of 192 bits is used, you need to append 0 after 192 bits of data and write 256 bits of data.
\end{tiplisting}

\subsubsection{Reading Data}

Reading data means reading data from the starting address of an ECDSA\_DS memory block and using this data as the output from the ECDSA\_DS. To be specific, reading data from an ECDSA\_DS memory block means reading $D[n][31:0]$ from the ``starting address of this ECDSA\_DS memory block + $4 \times n$".
For a 256-bit long data example:
\begin{itemize}
    \item read $D[0]$ from ``starting address"
    \item read $D[1]$ from ``starting address + 4"
    \item $\cdots$
    \item read $D[7]$ from ``starting address + 28"
\end{itemize}

\begin{tiplisting}
When the data size of 192 bits is used, only use the low 192 bits (6 blocks) of data.
\end{tiplisting}

\subsubsection{Padding the Message}\label{sec:ecdsa-msg-padding}

The SHA accelerator can only process message blocks of 512 bits. Thus, all the messages should be padded to a multiple of 512 bits before the hash operation.

Suppose that the length of the message \begin{math}M\end{math} is \begin{math} L_M \end{math} bits. Then \begin{math} M \end{math} shall be padded as introduced below:

\begin{enumerate}
    \item First, append the bit ``1" to the end of the message;
    \item Second, append \begin{math} L_A \end{math} bits of zeros, where \begin{math} L_A \end{math} is the smallest, non-negative solution to the equation \begin{math}L_M + 1 + L_A \ {\equiv} \ 448 \ mod \end{math} 512;
    \item Last, append the 64-bit block of value equal to the number \begin{math} L_M \end{math} expressed using a binary representation.
\end{enumerate}

For more details, please refer to \href{https://doi.org/10.6028/NIST.FIPS.180-4}{FIPS PUB 180-4 Spec} > Section ``Padding the Message".

\subsubsection{Parsing the Message}\label{sec:ecdsa-msg-parsing}

The message and its padding must be parsed into \begin{math}N\end{math} 512-bit message blocks: \begin{math}M^{(1)}\end{math}, \begin{math}M^{(2)}\end{math}, \dots, \begin{math}M^{(N)}\end{math}.

\begin{tiplisting}
\begin{enumerate}
\vspace{-2.5em}
    \item For more details about ``parsing the message", please refer to \href{https://doi.org/10.6028/NIST.FIPS.180-4}{FIPS PUB 180-4 Spec} > Section ``Parsing the Message".
    \item For more information on ``message block", please refer to \href{https://doi.org/10.6028/NIST.FIPS.180-4}{FIPS PUB 180-4 Spec} > Section ``Glossary of Terms and Acronyms".
\end{enumerate}
\end{tiplisting}


%%%%%%%%%%%%%%%%%%%%%%%%%
%%% Security Features %%%
%%%%%%%%%%%%%%%%%%%%%%%%%

\subsection{Security Features}
\label{sec:ecdsa-security_features}

To ensure the security of the operation process, the ECDSA\_DS module implements a variety of security functions.

\subsubsection{High Anti-Attack Performance}

\chipseries{}'s ECDSA\_DS module leverages ECC's enhanced anti-attack performance (refer to Section \ref{sec:ecc-anti-attack} \textit{\nameref{sec:ecc-anti-attack}}) every time it performs an operation. This means that each time a signature is generated and verified by ECDSA\_DS:

\begin{itemize}
    \item The execution latency is constant for a given operation.
    \item Power consumption variations are minimized.
\end{itemize}

This provides \chipseries{}'s ECDSA\_DS module strong anti-attack performance.

\subsubsection{State-Dependent Register Access Control}

\chipseries{}'s ECDSA\_DS module has implemented a state-dependent register access control mechanism to prevent any possibility of key theft by tampering with the configuration or accessing the data during the operation.

By implementing the state-dependent register access control, the accesses for ECDSA\_DS registers are designed to vary in different states. For example, \hyperref[regdesc:ECDSACONFREG]{ECDSA\_CONF\_REG} is only available for reading and writing when ECDSA\_DS is in the IDLE state. In this way, the configuration information is protected from reading or writing when ECDSA\_DS is in other states, such as LOAD, GAIN, and BUSY. For details about all ECDSA\_DS working states, please refer to Table \ref{tab:ecdsa-working-state}.

For detailed information on the state-dependent access control of each ECDSA\_DS register, please refer to Section \hyperref[sec:ecdsa-reg-summ]{Register Summary}.

\subsubsection{Hardware Occupation}

During the operation, the following hardware modules will be occupied by \chipname{}'s ECDSA\_DS:
\begin{itemize}
    \item SHA Accelerator
    \item ECC Accelerator
\end{itemize}
Among them, the SHA accelerator will be released when \hyperref[int:ECDSASHARELEASEINT]{ECDSA\_SHA\_RELEASE\_INT} is triggered, while the ECC accelerator will be occupied during the whole ECDSA\_DS operation.

\begin{tiplisting}
Hardware occupation is a mechanism to protect multiplexed modules and storage space. When a module is hardware occupied, the user will fail to:
\vspace{-0.5em}
\begin{itemize}
    \item read or write data to the module's registers or memories.
    \item disable the module clock.
    \item reset the module.
\end{itemize}
\vspace{-0.5em}
After hardware occupation is finished, the occupied module will be automatically reset. In addition, when the user performs a software reset to the master module, all the occupied modules will be reset at the same time.
\end{tiplisting}


%%%%%%%%%%%%%%%%%%%%%%%%%%%%%%
%%% Programming Procedures %%%
%%%%%%%%%%%%%%%%%%%%%%%%%%%%%%

\section{Programming Procedures}
\label{sec:ecdsa-prog-procedures}

\subsection{ECDSA\_DS Process}
The overall ECDSA\_DS process consists of the following six stages.

\begin{figure}[H]
\includegraphics[width=0.5\textwidth]{\modulefiles/figures/ecdsa-stage-changes.pdf}
\centering
\caption{ECDSA\_DS Process}
\end{figure}

The detailed programming procedures of each stage are described in the following sections.

\subsubsection{IDLE Stage}
In the IDLE stage:
\begin{enumerate}
	\item Configure the static parameters including eFuse bits:
	      \begin{enumerate}
		      \item ECDSA\_KEY: The value of private key $d$ in ECDSA\_DS. To correctly configure the key value in eFuse, users need to write the key value in KEY\regindex{n} (n = 0-5), and set the corresponding EFUSE\_KEY\_PURPOSE\_\regindex{n} as ECDSA\_KEY. Please refer to Chapter \ref{mod:efuse} \textit{\nameref{mod:efuse}} for more detailed configuration steps.
	      \end{enumerate}
	\item Configure \hyperref[regdesc:ECDSACONFREG]{ECDSA\_CONF\_REG}, including the following fields:
	      \begin{enumerate}
		      \item \hyperref[fielddesc:ECDSAWORKMODE]{ECDSA\_WORK\_MODE}: Select the working mode of ECDSA\_DS.
		      \item \hyperref[fielddesc:ECDSAECCCURVE]{ECDSA\_ECC\_CURVE}: Select the elliptic curve of ECDSA\_DS.
		      \item \hyperref[fielddesc:ECDSASOFTWARESETZ]{ECDSA\_SOFTWARE\_SET\_Z}: Configure whether to use direct input $z$.
		            \item \hyperref[fielddesc:ECDSADETERMINISTICK]{ECDSA\_DETERMINISTIC\_K}: Configure whether to use deterministic k.
		            \item \hyperref[fielddesc:ECDSADETERMINISTICLOOP]{ECDSA\_DETERMINISTIC\_LOOP}: Configure the loop number of deterministic k.
	      \end{enumerate}
	\item Configure the register field \hyperref[fielddesc:ECDSASTART]{ECDSA\_START} to enter the PREP stage.
\end{enumerate}
\begin{tiplisting}
For more details about Deterministic ECDSA, please refer to \href{https://www.rfc-editor.org/rfc/rfc6979}{RFC6979}.
\end{tiplisting}

\subsubsection{PREP Stage}

In the PREP stage, \hyperref[regdesc:ECDSASTATEREG]{ECDSA\_STATE\_REG} is BUSY, and ECDSA\_DS performs preparation.

Wait till the PREP stage to end by polling \hyperref[fielddesc:ECDSABUSY]{ECDSA\_BUSY} until it is not BUSY. Then ECDSA\_DS will automatically enter the LOAD stage.

\subsubsection{LOAD Stage}
In the LOAD stage:

\begin{enumerate}
	\item Provide input $z$ into ECDSA\_DS using one of the following options:
	      \begin{itemize}
		      \item Direct input $z$: write $z$ to \hyperref[tab:ecdsa-memory]{ECDSA\_MEM\_Z}.
		      \item ECDSA\_DS SHA interface: generate $z$ from the message. For details, please refer to Section \ref{sec:ecdsa-sha-interface}.
	      \end{itemize}
	\item According to the selected \hyperref[fielddesc:ECDSAWORKMODE]{ECDSA\_WORK\_MODE}, configure as follows:
	      \begin{itemize}
		      \item Signature Verification:
		            \begin{enumerate}
			            \item write signature ($r$, $s$) to \hyperref[tab:ecdsa-memory]{ECDSA\_MEM\_R} and \hyperref[tab:ecdsa-memory]{ECDSA\_MEM\_S}.
			            \item write public key ($Qx$, $Qy$) to \hyperref[tab:ecdsa-memory]{ECDSA\_MEM\_Qx} and \hyperref[tab:ecdsa-memory]{ECDSA\_MEM\_Qy}.
		            \end{enumerate}
		      \item Signature Generation:
		            \begin{enumerate}
			            \item no other configuration is required.
		            \end{enumerate}
                \item Public Key Export:
		            \begin{enumerate}
                            \item no other configuration is required.
		            \end{enumerate}
	      \end{itemize}
	\item Write 1 to \hyperref[fielddesc:ECDSALOADDONE]{ECDSA\_LOAD\_DONE}, indicating that the configuration is done. Then ECDSA\_DS will automatically enter the PROC stage.
\end{enumerate}

\subsubsection{PROC Stage}

In the PROC stage, \hyperref[regdesc:ECDSASTATEREG]{ECDSA\_STATE\_REG} is BUSY, and ECDSA\_DS performs operation based on the selected working mode.

Wait till the PROC stage to end by polling \hyperref[fielddesc:ECDSABUSY]{ECDSA\_BUSY} until it is not BUSY. Then ECDSA\_DS will automatically enter the either the GAIN stage or the POST stage depending on the selected working mode:
\begin{itemize}
    \item Signature generation: GAIN stage
    \item Signature verification: POST stage
\end{itemize}

\subsubsection{GAIN Stage}

When the Signature Generation mode or the Public Key Export mode is selected, ECDSA\_DS enters the GAIN stage after PROC stage:

\begin{enumerate}
	\item Read data from the ECDSA\_DS memory:
	      \begin{itemize}
		      \item read signature ($r$, $s$) from \hyperref[tab:ecdsa-memory]{ECDSA\_MEM\_R} and \hyperref[tab:ecdsa-memory]{ECDSA\_MEM\_S} only in the Signature Generation mode.
		      \item read public key ($Qx$, $Qy$) from \hyperref[tab:ecdsa-memory]{ECDSA\_MEM\_Qx} and \hyperref[tab:ecdsa-memory]{ECDSA\_MEM\_Qy} both in the Signature Generation mode and the Public Key Export mode.
	      \end{itemize}
	\item Write 1 to \hyperref[fielddesc:ECDSAGETDONE]{ECDSA\_GET\_DONE}, indicating that the GAIN stage is done. Then ECDSA\_DS will automatically enter the POST stage.
\end{enumerate}

\subsubsection{POST Stage}

In the POST stage, \hyperref[regdesc:ECDSASTATEREG]{ECDSA\_STATE\_REG} is BUSY, and ECDSA\_DS performs wrap-up work of the operation.

Wait till the POST stage to end by polling \hyperref[fielddesc:ECDSABUSY]{ECDSA\_BUSY} until it is not BUSY. Then ECDSA\_DS will automatically return to the IDLE stage.

\subsubsection{ECDSA\_DS SHA Interface}
\label{sec:ecdsa-sha-interface}

\chipseries{}'s ECDSA\_DS module can automatically execute hash operation and generates $z$ based on a direct input message.

For message hash, ECDSA\_DS supports SHA algorithms SHA-224 (only valid when P-192 is selected as the elliptic curve) and SHA-256.

To use the ECDSA\_DS SHA interface, complete the following steps:

\begin{enumerate}
	\item Pad the message by following the steps described in Section \ref{sec:ecdsa-msg-padding}.
	\item Parse the message and its padding into message blocks. See details in Section \ref{sec:ecdsa-msg-parsing}.
	\item \label{others:ecdsa-current-block} Process the current message block.
	      \begin{itemize}
		      \item Write the current message block into \hyperref[tab:ecdsa-memory]{ECDSA\_MEM\_M}.
	      \end{itemize}
	\item Start the ECDSA\_DS SHA interface\hyperref[others:ecdsa-note-sha]{$^1$}.
	      \begin{itemize}
		      \item If this is the first time to execute this step, write 1 to \hyperref[fielddesc:ECDSASHASTART]{ECDSA\_SHA\_START} to start the ECDSA\_DS SHA interface;
		      \item If this is not the first time to execute this step\hyperref[others:ecdsa-note-sha]{$^2$}, write 1 to \hyperref[fielddesc:ECDSASHACONTINUE]{ECDSA\_SHA\_CONTINUE} to continue the operation.
	      \end{itemize}
	\item Check the progress of the current message block processing by polling.
	      \begin{itemize}
		      \item Poll register \hyperref[fielddesc:ECDSASHABUSY]{ECDSA\_SHA\_BUSY} until it's IDLE, indicating that the interface has completed the operation for the current message block.
	      \end{itemize}

	\item Process the next message block:
	      \begin{itemize}
		      \item If there is message block to process, go back to Step \ref{others:ecdsa-current-block}.
		      \item If there is no more message blocks, exit.
	      \end{itemize}

\end{enumerate}

\begin{tiplisting}\label{others:ecdsa-note-sha}
\vspace{-2.5em}
\begin{enumerate}
\item In this step, the software can also write the next message block to be processed, if any, in register \hyperref[tab:ecdsa-memory]{ECDSA\_MEM\_M}, while the interface starts SHA operation, to save time.
\item You are resuming the ECDSA\_DS SHA interface with the previously paused operation.
\end{enumerate}
\end{tiplisting}

\subsection{Clocks and Resets}

\chipname{}'s ECDSA\_DS module has one clock module and one reset module. It is activated by setting the \hyperref[fielddesc:PCRECDSACLKEN]{PCR\_ECDSA\_CLK\_EN} bit and clearing the \hyperref[fielddesc:PCRECDSARSTEN]{PCR\_ECDSA\_RST\_EN} bit. For details on how to configure the ECDSA\_DS clock and reset, please refer to Chapter \ref{mod:resclk} \textit{\nameref{mod:resclk}}.

\subsection{Interrupts}
\label{sec:ecdsa-interrupt}

\chipname{}'s ECDSA\_DS module can generate the interrupt signal ECDSA\_INTR and send it to \hyperref[mod:intmtrx]{Interrupt Matrix}.

The ECDSA\_DS module has four interrupt sources that can generate the ECDSA\_INTR interrupt signal as follows. Each interrupt source can be configured by a common set of registers that are described in Section \hyperref[interrupt-config-registers]{\textit{Interrupt Configuration Registers}}.

\begin{itemize}
	\item \label{int:ECDSAPREPDONEINT}ECDSA\_PREP\_DONE\_INT: triggered on the completion of the PREP stage. This interrupt source is configured by the following registers:
	      \begin{itemize}
		      \item \hyperref[fielddesc:ECDSAPREPDONEINTRAW]{ECDSA\_PREP\_DONE\_INT\_RAW}: stores the raw interrupt status of ECDSA\_PREP\_DONE\_INT.
		      \item \hyperref[fielddesc:ECDSAPREPDONEINTST]{ECDSA\_PREP\_DONE\_INT\_ST}: indicates the status of the ECDSA\_PREP\_DONE\_INT interrupt. This field is generated by enabling/disabling the \hyperref[fielddesc:ECDSAPREPDONEINTRAW]{ECDSA\_PREP\_DONE\_INT\_RAW} field.
		      \item \hyperref[fielddesc:ECDSAPREPDONEINTENA]{ECDSA\_PREP\_DONE\_INT\_ENA}: enables/disables the ECDSA\_PREP\_DONE\_INT interrupt.
		      \item \hyperref[fielddesc:ECDSAPREPDONEINTCLR]{ECDSA\_PREP\_DONE\_INT\_CLR}: set this bit to clear the ECDSA\_PREP\_DONE\_INT interrupt status. By setting this bit to 1, fields \hyperref[fielddesc:ECDSAPREPDONEINTRAW]{ECDSA\_PREP\_DONE\_INT\_RAW} and  \hyperref[fielddesc:ECDSAPREPDONEINTST]{ECDSA\_PREP\_DONE\_INT\_ST} will be cleared.
	      \end{itemize}
	\item \label{int:ECDSAPROCDONEINT}ECDSA\_PROC\_DONE\_INT: triggered on the completion of the PROC stage. This interrupt source is configured by the following registers:
	      \begin{itemize}
		      \item \hyperref[fielddesc:ECDSAPROCDONEINTRAW]{ECDSA\_PROC\_DONE\_INT\_RAW}: stores the raw interrupt status of ECDSA\_PROC\_DONE\_INT.
		      \item \hyperref[fielddesc:ECDSAPROCDONEINTST]{ECDSA\_PROC\_DONE\_INT\_ST}: indicates the status of the ECDSA\_PROC\_DONE\_INT interrupt. This field is generated by enabling/disabling the \hyperref[fielddesc:ECDSAPROCDONEINTRAW]{ECDSA\_PROC\_DONE\_INT\_RAW} field.
		      \item \hyperref[fielddesc:ECDSAPROCDONEINTENA]{ECDSA\_PROC\_DONE\_INT\_ENA}: enables/disables the ECDSA\_PROC\_DONE\_INT interrupt.
		      \item \hyperref[fielddesc:ECDSAPROCDONEINTCLR]{ECDSA\_PROC\_DONE\_INT\_CLR}: set this bit to clear the ECDSA\_PROC\_DONE\_INT interrupt status. By setting this bit to 1, fields \hyperref[fielddesc:ECDSAPROCDONEINTRAW]{ECDSA\_PROC\_DONE\_INT\_RAW} and  \hyperref[fielddesc:ECDSAPROCDONEINTST]{ECDSA\_PROC\_DONE\_INT\_ST} will be cleared.
	      \end{itemize}
	\item \label{int:ECDSAPOSTDONEINT}ECDSA\_POST\_DONE\_INT: triggered on the completion of the POST stage. This interrupt source is configured by the following registers:
	      \begin{itemize}
		      \item \hyperref[fielddesc:ECDSAPOSTDONEINTRAW]{ECDSA\_POST\_DONE\_INT\_RAW}: stores the raw interrupt status of ECDSA\_POST\_DONE\_INT.
		      \item \hyperref[fielddesc:ECDSAPOSTDONEINTST]{ECDSA\_POST\_DONE\_INT\_ST}: indicates the status of the ECDSA\_POST\_DONE\_INT interrupt. This field is generated by enabling/disabling the \hyperref[fielddesc:ECDSAPOSTDONEINTRAW]{ECDSA\_POST\_DONE\_INT\_RAW} field.
		      \item \hyperref[fielddesc:ECDSAPOSTDONEINTENA]{ECDSA\_POST\_DONE\_INT\_ENA}: enables/disables the ECDSA\_POST\_DONE\_INT interrupt.
		      \item \hyperref[fielddesc:ECDSAPOSTDONEINTCLR]{ECDSA\_POST\_DONE\_INT\_CLR}: set this bit to clear the ECDSA\_POST\_DONE\_INT interrupt status. By setting this bit to 1, fields \hyperref[fielddesc:ECDSAPOSTDONEINTRAW]{ECDSA\_POST\_DONE\_INT\_RAW} and  \hyperref[fielddesc:ECDSAPOSTDONEINTST]{ECDSA\_POST\_DONE\_INT\_ST} will be cleared.
	      \end{itemize}
	\item \label{int:ECDSASHARELEASEINT}ECDSA\_SHA\_RELEASE\_INT: triggered when SHA is released.This interrupt source is configured by the following registers:
	      \begin{itemize}
		      \item \hyperref[fielddesc:ECDSASHARELEASEINTRAW]{ECDSA\_SHA\_RELEASE\_INT\_RAW}: stores the raw interrupt status of ECDSA\_SHA\_RELEASE\_INT.
		      \item \hyperref[fielddesc:ECDSASHARELEASEINTST]{ECDSA\_SHA\_RELEASE\_INT\_ST}: indicates the status of the ECDSA\_SHA\_RELEASE\_INT interrupt. This field is generated by enabling/disabling the \hyperref[fielddesc:ECDSASHARELEASEINTRAW]{ECDSA\_SHA\_RELEASE\_INT\_RAW} field via \hyperref[fielddesc:ECDSASHARELEASEINTENA]{ECDSA\_SHA\_RELEASE\_INT\_ENA}.
		      \item \hyperref[fielddesc:ECDSASHARELEASEINTENA]{ECDSA\_SHA\_RELEASE\_INT\_ENA}: enables/disables the ECDSA\_SHA\_RELEASE\_INT interrupt.
		      \item \hyperref[fielddesc:ECDSASHARELEASEINTCLR]{ECDSA\_SHA\_RELEASE\_INT\_CLR}: set this bit to clear the ECDSA\_SHA\_RELEASE\_INT interrupt status. By setting this bit to 1, fields \hyperref[fielddesc:ECDSASHARELEASEINTRAW]{ECDSA\_SHA\_RELEASE\_INT\_RAW} and  \hyperref[fielddesc:ECDSASHARELEASEINTST]{ECDSA\_SHA\_RELEASE\_INT\_ST} will be cleared.
	      \end{itemize}
\end{itemize}

\begin{tiplisting}
For definitions of \textit{interrupt}, \textit{interrupt signal}, \textit{interrupt source}, and their correlations, please refer to Chapter \ref{mod:intmtrx} \textit{\nameref{mod:intmtrx}} > Section \ref{sec:intmtx-int-terminology} \textit{\nameref{sec:intmtx-int-terminology}}.
\end{tiplisting}

\newpage
\section{Memory Blocks}

ECDSA\_DS's memory blocks store input data and output data of the operation.

\begin{table}[H]
    \centering
    \caption{ECDSA\_DS Memory Blocks} \label{tab:ecdsa-memory}
    \begin{threeparttable}
    \begin{tabular}{ | l | l | l | l | l | l | }
    \hline
        \rowcolor{lightgray}
        \textbf{Memory} & \textbf{Size (byte)}  &  \textbf{Starting Address}\tnote{*} &  \textbf{Ending Address} \tnote{*} & \textbf{Access}   \\ \hline
        ECDSA\_MEM\_M    & 64  & 0x280 & 0x2BF & R/W \\ \hline
        ECDSA\_MEM\_R    & 32  & 0xA00 & 0xA1F & R/W \\ \hline
        ECDSA\_MEM\_S    & 32  & 0xA20 & 0xA3F & R/W \\ \hline
        ECDSA\_MEM\_Z    & 32  & 0xA40 & 0xA5F & R/W \\ \hline
        ECDSA\_MEM\_Qx   & 32  & 0xA60 & 0xA7F & R/W \\ \hline
        ECDSA\_MEM\_Qy   & 32  & 0xA80 & 0xA9F & R/W \\ \hline
    \end{tabular}
    \begin{tablenotes}
        \item[*] Address offset related to ECDSA\_DS base address is provided in Table \ref{tab:sysmem-base-address} in Chapter \ref{mod:sysmem} \textit{\nameref{mod:sysmem}}.
    \end{tablenotes}
    \end{threeparttable}
\end{table}


% >>> If your module has no registers, delete sample text
%       for Sections Register Summary and Registers below <<<

%%%%%%%%%%%%%%%%%%%%%%%%
%%% Register Summary %%%
%%%%%%%%%%%%%%%%%%%%%%%%

% >>> Update the sample text below and insert
%     the info on register summary into the subfile <<<

\newpage
\hypertarget{ecdsa-reg-summ}{}
\section{Register Summary}
\label{sec:ecdsa-reg-summ}

The addresses in this section are relative to ECDSA\_DS base address provided in Table \ref{tab:sysmem-base-address} in Chapter \ref{mod:sysmem} \textit{\nameref{mod:sysmem}}.

To enhance security, ECDSA\_DS registers have different read and write permissions in different operating states. The following is the abbreviation and corresponding relationship of each state:

\begin{itemize}
    \item PI: IDLE state, corresponding to IDLE stage.
    \item PL: LOAD state, corresponding to LOAD stage.
    \item PG: GAIN state, corresponding to GAIN stage.
    \item PB: BUSY state, corresponding to PREP, PROC, and POST stage.
\end{itemize}

Other abbreviations given in Column \textbf{Access} are explained in Section \hyperref[glossary-access-types]{Access Types for Registers}.

\subfile{\modulefiles/23-ECDSA-reg-summ__EN}


%%%%%%%%%%%%%%%%%
%%% Registers %%%
%%%%%%%%%%%%%%%%%

% >>> Update the sample text below and insert
%      the info on registers into the subfile <<<

\newpage
\section{Registers}

The addresses in this section are relative to ECDSA\_DS base address provided in Table \ref{tab:sysmem-base-address} in Chapter \ref{mod:sysmem} \textit{\nameref{mod:sysmem}}.

For how to program reserved fields, please refer to Section \hyperref[programming-reserved-field]{\textit{Programming Reserved Register Field}}.

\subfile{\modulefiles/23-ECDSA-reg__EN}


\end{document}
